\section{Discussion} \label{sec:diss}

\subsection{Implications for Package Security}

Source-only package analysis is incomplete when distributions contain executable bytecode without matching source.  Package scanners should inventory \texttt{.pyc} files, inspect marshalled code objects, flag bytecode/source mismatches, and preserve decompilation failures as evidence for manual review.  A failed source recovery step should not be treated as a clean result.

\subsection{Implications for Runtime Hardening}

CPython need not support arbitrary malformed bytecode as trusted input, but it should fail predictably at documented validation boundaries.  Assertions, aborts, sanitizer findings, and segmentation faults indicate hardening opportunities even when the triggering artifact is malformed.  The source-reproduction analysis helps separate findings reproduced by selected source workflows from findings that, in this study, require direct bytecode construction.

\subsection{Implications for Malware Analysis}

Compiled Python payloads complicate triage because the analyst may see a small source-level loader while the payload resides in a code object.  Bytecode-aware signatures, opcode-level behavior extraction, and explicit handling for version-specific bytecode can improve malware analysis workflows.  Decompilation remains useful, but it should be one component of a broader bytecode analysis pipeline.

\subsection{Threats to Validity}

Dataset bias is a primary threat: PyPI packages may not represent all Python deployments, including containers, application bundles, or private release artifacts.  Tool bias is also possible because analyzability and source-reproduction results depend on selected tools and versions.  Runtime fuzzing may overproduce malformed inputs that ordinary source would not generate, and finding deduplication can over- or under-count unique failures.  We mitigate these threats by reporting failure categories explicitly, using version-matched interpreters, and separating bytecode robustness from tool-bounded source reproduction.

\subsection{Limitations}

Several limitations remain.  First, {\tech} cannot prove that all dynamically loaded bytecode has been recovered; loaders may decrypt, decompress, download, or synthesize code objects in ways that are not visible statically.  Second, the PyPI dataset does not cover every Python deployment format.  Third, bytecode transparency based on nearby source is conservative and does not by itself establish semantic equivalence.  Fourth, decompiler results reflect the selected tools and versions; future tools may recover artifacts that are currently classified as tool failures.  Fifth, robustness findings from malformed bytecode are not automatically security vulnerabilities.  They require manual triage, source-reproduction analysis, and responsible disclosure when appropriate.
